%------------------------------------------------------------------------------------
%  SECTION 6.1 : GÉOMÉTRIES MULTIPLES À LA VOLÉE
%------------------------------------------------------------------------------------
\section{Alternance Dynamique des Géométries de Page}
\label{sec:geometries_multiples}

Dans un mémoire ou un rapport d'ingénierie, une taille de marge unique pour l'ensemble du volume est inadaptée :
la page de couverture et les grands tableaux requièrent des marges étroites pour maximiser l'espace utile, le corps du texte nécessite une largeur de colonne ergonomique (65 à 75 caractères par ligne pour un confort visuel optimal), tandis que la phase de relecture collaborative exige une marge latérale généreuse pour accueillir les annotations.

En \LaTeX{} standard, changer de géométrie en cours de document provoque souvent des décalages imprévisibles dans les en-têtes ou rompt l'alignement recto-verso.
\texttt{ClasseGuillaume} résout cette difficulté en encapsulant la gestion du paquet \textsf{geometry} et du moteur \textsf{typearea} de KOMA-Script au sein de trois commandes directes et sécurisées.

\subsection{Les trois commandes de géométrie de la classe}
\label{subsec:trois_commandes_geometrie}

\begin{itemize}
  \item \textbf{\texttt{\textbackslash geometrieWord} :}
  Fixe des marges symétriques de \qty{2.5}{\centi\meter} sur les quatre bords et alloue une marge latérale dédiée de \qty{2.2}{\centi\meter} (\texttt{marginparwidth}) avec un séparateur de \qty{3}{\milli\meter} (\texttt{marginparsep}).
  C'est la géométrie par excellence pour le corps de texte rédigé et la phase de relecture collaborative avec les notes \texttt{\textbackslash todoH}.
  
  \item \textbf{\texttt{\textbackslash geometrieLarge} (ou \texttt{\textbackslash geometrieExemple}) :}
  Applique des marges compactes de \qty{1.2}{\centi\meter} à gauche et à droite, et \qty{1.6}{\centi\meter} en haut et en bas.
  Elle annule l'espace marginal pour offrir une largeur utile maximale de \qty{18.6}{\centi\meter}, idéale pour la page de garde, les posters de synthèse, ou les grands tableaux récapitulatifs.

  \item \textbf{\texttt{\textbackslash geometrieStandard} :}
  Restaure instantanément la mise en page canonique calculée par l'algorithme typographique de KOMA-Script (\texttt{DIV = 12}), garantissant des proportions harmonieuses selon les canons du livre d'art.
\end{itemize}

\begin{boiteSynthese}{Règle de bascule propre}
  Chacune de ces commandes intègre automatiquement une directive \texttt{\textbackslash cleardoublepage} avant d'appliquer les nouvelles marges.
  Cette précaution fondamentale garantit qu'une nouvelle géométrie commence toujours sur une page neuve et prévient tout débordement de texte entre deux pages consécutives aux proportions divergentes.
\end{boiteSynthese}

